\documentclass[german, parskip=full]{scrartcl}
\usepackage[utf8]{inputenc}  % Kodierung der Datei
\usepackage[ngerman]{babel}  % Sprache auf Deutsch 
\usepackage{color}
\usepackage{graphicx, gensymb}
\usepackage{amsfonts, cancel, amssymb}
\usepackage{amsmath, amsthm, array, esvect, esint}
\usepackage{booktabs, multirow}  % schönere Tabellen
\usepackage{caption}  % Anpassung der Captions
\usepackage{scrlayer-scrpage}    % Manipulation des Seitenstils
\pagestyle{scrheadings}  % Stil für die Seite setzen
\automark{section}  % Abschnittsnamen als Seitenbeschriftung 
\setheadsepline{.5pt}    % Kopzeile durch Linie abtrennen
\usepackage[hidelinks]{hyperref}  % Links und weitere PDF-Features
\usepackage{typearea, tikz, pgffor, verbatim}
\usetikzlibrary{calc, quotes}
\usepackage[cache=false, outputdir=build]{minted}

\definecolor{bg}{rgb}{0.9,0.9,0.9}
\newcommand\numberthis{\addtocounter{equation}{1}\tag{\theequation}}
\newcommand{\bra}{}
\newcommand{\kom}{}
\newcommand{\brak}{}
\newcommand{\braket}{}
\newcommand{\intx}{}
\newcommand{\intr}{}
\newcommand{\kla}{}
\newcommand{\klam}{}
\newcommand{\klammer}{}
\newcommand{\oper}{}
\newcommand{\tecxt}{}
\newcommand{\Bra}{}
\newcommand{\Ket}{}
\newcommand{\I}{\text{i}}
\newcommand{\e}{\text{e}}
\newcommand*{\Fourier}{\textsc{Fourier}\ }
\newcommand*{\F}{\mathcal{F}}
\newcommand*{\sinc}{\text{sinc\,}}
\newcommand{\conv}{\mathop{\scalebox{3.5}{\raisebox{-0.38ex}{\makebox[0.5\width][c]{$\ast$}}}}\limits}

\newcommand*{\titel}{Klausur}
\newcommand*{\autor}{Joachim Schwardt}

\hypersetup{pdfauthor={\autor}, pdftitle={\titel}}

%\titlehead{titlehead}
\subject{Quantentheorie 2}
\title{\titel}
\author{100 Punkte entsprechen 100\%}
\date{am 05.08.2021}

\begin{document}

%\begin{titlepage}
\maketitle  % Titel setzen
%\tableofcontents  % Inhaltsverzeichnis setzen
%\end{titlepage}

\section{Diracgleichung freies Teilchen (20 Punkte)}
Betrachten Sie einen 4-Impuls 
\begin{align}
(p^\mu)&=(E,p,p,0)^\top\ \tecxt\text{mit}\ E^2=m^2+2p^2 \label{eqn:VierImpuls}
\end{align}
und einen Diracspinor der Form
\begin{align}
u(p)&=(1,\e^{I\frac{\pi}{4}},a,b)^\top .  \label{eqn:Diracspinor}
\end{align}
\begin{enumerate}
\item[a)] Werten Sie $p^\mu p_\mu$ aus! (2 Punkte)
\item[b)] Bestimmen Sie in dem Spinor $u(p)$ die Koeffizienten $a$ und $b$ so, dass $\psi(x^\mu) = u(p)\e^{-\I x^\mu p_\mu}$ die freie Diracgleichung löst. 
Nutzen Sie die Diracdarstellung der $\gamma^\mu $-Matrizen! (8 Punkte)
\item[c)] Ist der erhaltene Spinor ein Eigenspinor zum Spin-$z$-Operator $S_z$? 
Interpretieren/kommentieren Sie das Ergebnis! (3 Punkte)
\item[d)] Bei Drehungen im 3-dimensionalen Raum werden Spinoren mit gewissen Matrizen multi-
pliziert (s. Anhang mit Formelsammlung). 
Um welche Achse ($x$-, $y$-, oder $z$-Achse) und um
welchen Winkel müssen Sie den Spinor drehen, um einen Eigenspinor zum Spin-$x$-Operator
$S_x$ zu erhalten? 
Zeigen Sie die expliziten Rechenschritte! (7 Punkte)
\end{enumerate}
\newpage
\subsubsection*{Lösungsvorschlag}
\begin{enumerate}
\item[a)] Wie zu erwarten gilt $p^\mu p_\mu = (p^0)^2-(p^1)^2-(p^2)^2-(p^3)^2=E^2-2p^2=m^2$. 
\item[b)] Die Diracgleichung lautet hier
\begin{align*}
0&=(\I\gamma^\mu\partial_\mu -m)\psi=\e^{-\I x^\mu p_\mu}(\gamma^\mu p_\mu -m)u(p)\propto (E\gamma^0-p(\gamma^1+\gamma^2)-m)u(p).
\end{align*}
Mit dem gegebenen Spinor findet man das Gleichungssystem
\begin{align*}
0&=\begin{pmatrix}
E-m&0&0&-p(1-\I)\\
0&E-m&-p(1+\I)&0\\
0&p(1-\I)&-(E+m)&0\\
p(1+\I)&0&0&-(E+m)
\end{pmatrix}\begin{pmatrix}
1\\ \e^{\I\frac{\pi}{4}}\\ a\\ b
\end{pmatrix}=\begin{pmatrix}
(E-m)-pb(1-\I)\\ (E-m)\e^{\I\frac{\pi}{4}}-pa(1+\I)\\ p(1-\I)\e^{\I\frac{\pi}{4}}-a(E+m)\\ p(1+\I) -b(E+m)
\end{pmatrix}.
\end{align*}
Kombinieren der ersten zwei Gleichungen liefert
\begin{align*}
pb(1-\I)&=pa(1+\I)\e^{-\I\frac{\pi}{4}}&&\Rightarrow &b&=a\e^{\I\frac{\pi}{4}}.
\end{align*}
Analog erhält man aus dem zweiten Paar die äquivalente Bedingung
\begin{align*}
pa(1+\I)&=pb(1-\I)\e^{\I\frac{\pi}{4}}&&\Rightarrow &a&=b\e^{-\I\frac{\pi}{4}}.
\end{align*}
Einsetzen in die erste Gleichung liefert
\begin{align*}
E-m&=pb(1-\I)=pa\sqrt{2}\e^{-\I\frac{\pi}{4}}\e^{\I\frac{\pi}{4}}=pa\sqrt{2}&&\Rightarrow &a&=\frac{E-m}{\sqrt{2}p}.
\end{align*}
Über die vorherige Relation erhält man auch
\begin{align*}
b&=a\e^{\I\frac{\pi}{4}}=\frac{E-m}{p}(1+\I).
\end{align*}
Definiere $\kappa:=\frac{E-m}{\sqrt{2}p}$.
Dann lautet der Eigenspinor
\begin{align*}
u(p)&=(1,\e^{\I\frac{\pi}{4}},\kappa,\kappa\e^{\I\frac{\pi}{4}})^\top .
\end{align*}
\item[c)] Dieser Spinor ist kein Eigenspinor zu $S_z$.
Betrachte dazu das Eigenwertproblem
\begin{align*}
S_zu(p)&=\frac{1}{2}\begin{pmatrix}
1&0&0&0\\
0&-1&0&0\\
0&0&1&0\\
0&0&0&-1
\end{pmatrix}\begin{pmatrix}
1\\ \e^{\I\frac{\pi}{4}} \\ \kappa \\ \kappa \e^{\I\frac{\pi}{4}}
\end{pmatrix}=\frac{1}{2}\begin{pmatrix}
1\\ -\e^{\I\frac{\pi}{4}} \\ \kappa \\ -\kappa \e^{\I\frac{\pi}{4}}
\end{pmatrix}\overset{!}{=}\lambda u(p).
\end{align*}
Offensichtlich sind die erste und dritte Gleichung durch $\lambda=+\frac{1}{2}$ erfüllt.
Allerdings ergibt sich aus der zweiten und vierten Gleichung $\lambda=-\frac{1}{2}$,  woraus sich ein Widerspruch ergibt.
Folglich ist $u(p)$ kein Eigenspinor zu $S_x$.
\item[d)] Es ist sinnvoll, zunächst die allgemeine Eigenwertgleichung für $S_x$ zu betrachten.
Mit $u=(u_1,u_2,u_3,u_4)^\top$ lautet diese
\begin{align*}
S_xu&=\frac{1}{2}\begin{pmatrix}
0&1&0&0\\
1&0&0&0\\
0&0&0&1\\
0&0&1&0
\end{pmatrix}\begin{pmatrix}
u_1\\ u_2 \\ u_3 \\ u_4
\end{pmatrix}=\frac{1}{2}\begin{pmatrix}
u_2 \\ u_1 \\ u_4 \\ u_3
\end{pmatrix}\overset{!}{=}\lambda u.
\end{align*}
Aus $u_2=2\lambda u_1$ und $u_1=2\lambda u_2$ folgt dann $1=4\lambda^2$, bzw. $\lambda=\pm\frac{1}{2}$.
Damit lautet die allgemeine Form eines $S_x$-Eigenspinors
\begin{align*}
u=(u_1,\pm u_1,u_3,\pm u_3)^\top .
\end{align*}
Um diese Form für $u(p)$ zu erreichen, ist es ausreichend, den Spinor in die Form $(\e^{\I\frac{\pi}{8}},\e^{\I\frac{\pi}{8}},\kappa\e^{\I\frac{\pi}{8}},\kappa\e^{\I\frac{\pi}{8}})^\top$ zu überführen.
Dazu genügt eine Drehung um die $z$-Achse mit $\theta=\frac{\pi}{4}\equiv 45\degree$.
Betrachte zum Nachweis 
\begin{align*}
S\kla\left( {R_x\kla\left( {\theta=\frac{\pi}{4}} \right)} \right)u(p)&=\begin{pmatrix}
\e^{\I\frac{\pi}{8}}&0&0&0\\
0&\e^{-\I\frac{\pi}{8}}&0&0\\
0&0&\e^{\I\frac{\pi}{8}}&0\\
0&0&0&\e^{-\I\frac{\pi}{8}}
\end{pmatrix}\begin{pmatrix}
1\\ \e^{\I\frac{\pi}{4}} \\ \kappa \\ \kappa \e^{\I\frac{\pi}{4}}
\end{pmatrix}=\begin{pmatrix}
\e^{\I\frac{\pi}{8}} \\ \e^{\I\frac{\pi}{8}} \\ \kappa\e^{\I\frac{\pi}{8}} \\ \kappa\e^{\I\frac{\pi}{8}}
\end{pmatrix}.
\end{align*}
Dieser Eigenspinor hat den Eigenwert $\lambda=+\frac{1}{2}$.
\end{enumerate}
\newpage
\section{Diracgleichung mit Potential (20 Punkte)}
Betrachten Sie die Diracgleichung für ein Elektron mit Ladung $Q=-e$ und Wechselwirkung mit dem elektromagnetischen Feld (d.h. kovariante Ableitung $\I D_\mu =\I\partial_\mu +eA_\mu$ ).
Spezialisieren Sie auf den folgenden Fall einer zeitunabhängigen Potentialstufe 
\begin{align}
-e\vec{A}(x^\mu)&=0\\
-e A^0(x^\mu)&=V(z)=\begin{cases}0&\tecxt\text{für}\ z<0 \\ V_0&\tecxt\text{für}\ z\ge 0\end{cases}.
\end{align}
Für die Wellenfunktion machen wir den einfachen Ansatz $\psi (x^\mu)=f_{1,2}(z)\e^{-\I Et}$ mit
\begin{align}
z<0:\ \ f_1(z)&=\begin{pmatrix}
1\\ 0\\ a\\ 0
\end{pmatrix}\e^{\I pz}+\alpha\begin{pmatrix}
1\\ 0\\ -a\\ 0
\end{pmatrix}\e^{-\I pz},&&\tecxt\text{und}&z\ge 0:\ \ f_2(z)&=\beta\begin{pmatrix}
1\\ 0\\ a'\\ 0
\end{pmatrix}\e^{\I p'z}.\label{eqn:AnsatzPotentialstufe}
\end{align}
Ferner betrachten wir den konkreten Fall eines leicht relativistischen Teilchens und einer sehr
hohen Potentialstufe
\begin{align}
E&=\frac{5}{4}m&&\tecxt\text{und}&V_0&=E+\frac{13}{5}m.\label{eqn:KonkreteEnergieWerte}
\end{align}
\begin{enumerate}
\item[a)] Zeigen Sie zunächst, dass die Diracgleichung (mit Diracdarstellung der $\gamma^\mu$-Matrizen) für $z\ge 0$ mit dem entsprechenden Ansatz (\ref{eqn:AnsatzPotentialstufe}) gelöst wird -- bestimmen Sie die nötigen Parameter $p'>0$ und $a'$! 
Setzen Sie auch die Werte (\ref{eqn:KonkreteEnergieWerte}) ein! (6 Punkte)
\item[b)] Zeigen Sie dann, dass die Diracgleichung für $z < 0$ mit dem Ansatz (\ref{eqn:AnsatzPotentialstufe}) gelöst wird --
bestimmen Sie die nötigen Parameter $p>0$ und $a$!
Hierbei können Sie auch ohne Rechnung eine
einfache Ersetzungsregel auf das Ergebnis für $z \ge 0$ anwenden. 
Setzen Sie auch die Werte
(\ref{eqn:KonkreteEnergieWerte}) ein! (6 Punkte)
\item[c)] Bei $z = 0$ sollte die Anschlussbedingung $f_1 (0) = f_2 (0)$ gelten. 
Geben Sie einen kurzen
Kommentar, warum die Anschlussbedingung $f_1'(0)=f_2'(0)$ hier nicht gilt. 
Bestimmen Sie
damit den Reflexionskoeffizienten $R = |\alpha|^2$ für die Werte (\ref{eqn:KonkreteEnergieWerte})! (5 Punkte)
\item[d)] Geben Sie einen kurzen Kommentar zum Ergebnis. 
Welche Tatsache deutet sich durch
das Ergebnis an? (3 Punkte)
\end{enumerate}
\newpage
\subsubsection*{Lösungsvorschlag}
\begin{enumerate}
\item[a)] Mit $\psi = u\e^{\I pz}$ und $u=(1,0,a',0)^\top$ lautet hier die Diracgleichung
\begin{align*}
0&=(\I\gamma^\mu D_\mu -m)\psi\propto(p^\mu\partial_\mu + e\gamma^\mu A_\mu  -m)u=([E-V_0]\gamma^0-p'\gamma^3-m)u=\\
&=\begin{pmatrix}
E-V_0-m&0&-p'&0\\
0&E-V_0-m&0&p'\\
p'&0&V_0-E-m&0\\
0&-p'&0&V_0-E-m
\end{pmatrix}\begin{pmatrix}
1 \\ 0 \\ a' \\ 0
\end{pmatrix} .
\end{align*}
Daraus folgen die Gleichungen
\begin{align*}
p'a'&=E-V_0-m&&\tecxt\text{und}&p'&=a'(E-V_0+m),
\end{align*}
und somit schließlich
\begin{align*}
a'&=\sqrt{\frac{E-V_0-m}{E-V_0+m}}\overset{(\ref{eqn:KonkreteEnergieWerte})}{=}\sqrt{\frac{\frac{13}{5}-1}{\frac{13}{5}+1}}=\frac{2}{3},\\
p'&=\sqrt{(E-V_0-m)(E-V_0+m)}\overset{(\ref{eqn:KonkreteEnergieWerte})}{=}\sqrt{\kla\left( {\frac{13}{5}-1} \right)\kla\left( {\frac{13}{5}+1} \right)}m=\frac{12}{5}m.
\end{align*}
\item[b)] Für $z<0$ führt die Rechnung auf ein sehr ähnliches Erebnis.
Zunächst entällft $V_0$, man verwende also die Ersetzungsregel $E-V_0\rightarrow E$.
Im zweiten Term gilt dazu noch $a\rightarrow -a$ und $p\rightarrow -p$.
Diese ändern allerdings nichts an den zwei relevanten Gleichungen, da dort $p$ und $a$ nur als Produkt bzzw. Quotient vorkommen.
Folglich erhalten wir
\begin{align*}
a&=\sqrt{\frac{E-m}{E+m}}\overset{(\ref{eqn:KonkreteEnergieWerte})}{=}\sqrt{\frac{\frac{5}{4}-1}{\frac{5}{4}+1}}=\frac{1}{3},\\
p&=\sqrt{(E-m)(E+m)}\overset{(\ref{eqn:KonkreteEnergieWerte})}{=}\sqrt{\kla\left( {\frac{5}{4}-1} \right)\kla\left( {\frac{5}{4}+1} \right)}m=\frac{3}{4}m.
\end{align*}
\item[c)] Die Ableitungen können bei $z=0$ unstetig sein, da auch das Potential hier unstetig ist.
Aus $f_1(0)=f_2(0)$ folgt zunächst
\begin{align*}
1+\alpha&=\beta&&\tecxt\text{und}&a-\alpha a&=\beta a'.
\end{align*}
Einsetzen liefert dann
\begin{align*}
\frac{a}{a'}\kla\left( {1-\alpha} \right)&=1+\alpha\\
\Rightarrow\ \ \ \alpha &=\frac{\frac{a}{a'}-1}{\frac{a}{a'}+1}=\frac{a-a'}{a+a'}\overset{(\ref{eqn:KonkreteEnergieWerte})}{=}\frac{\frac{1}{3}-\frac{2}{3}}{\frac{1}{3}+\frac{2}{3}}=-\frac{1}{3}.
\end{align*}
Damit ist der Reflexionskoeffizient gegeben durch $R=|\alpha|^2=\frac{1}{9}$. 
\item[d)] Der Transmissionskoeffizient ist mit $T=1-R=\frac{8}{9}$ überraschend groß.
Im nichtrelativistischen Fall wird der größte Anteil einer Welle reflektiert, wenn die Energie signifikant kleiner als die Potentialstufe ist.
Anscheinend ist dies bereits für leicht relativistische Teilchen nicht mehr der Fall.
\end{enumerate}
\newpage
\section{Identische Teilchen (20 Punkte)}
Es seien folgende 1-Teilchenzustände bekannt.
Orthonormierte Eigenzustände $\Ket | {\alpha} \rangle$ für $\alpha\in\{a,b\}$ zu irgendeiner Observablen $\hat{O}$, sowie orthonormierte Energieeigenzustände $\Ket | {j} \rangle$ für $j\in\{1,2\}$.
Diese seien über
\begin{align}
\Ket | {1} \rangle&=c\Ket | {a} \rangle-s\Ket | {b} \rangle&&\tecxt\text{und}&\Ket | {2} \rangle&=s\Ket | {a} \rangle+c\Ket | {b} \rangle
\end{align}
miteinander verknüpft, wobei $c^2+s^2=1$.
\begin{enumerate}
\item[a)] Betrachten Sie zwei unterscheidbare Teilchen, präpariert im Zustand $|a,b\rangle$. 
Geben Sie die Wahrscheinlichkeitsamplituden an, bei einer Energiemessung die Endzustände
\begin{align}
\Ket | {1,1} \rangle,&&\Ket | {1,2} \rangle,&&\Ket | {2,1} \rangle,&&\Ket | {2,2} \rangle
\end{align}
zu finden. (4 Punkte)
\item[b)] Was ist die Summe der vier entsprechenden Wahrscheinlichkeiten? (2 Punkte)
\item[c)] Betrachten Sie zwei identische Bosonen, präpariert im auf 1 normierten Zustand $\Ket | {a,b} \rangle^+$. 
Geben Sie die Wahrscheinlichkeitsamplituden an, bei einer Energiemessung die Endzustände
\begin{align}
\Ket | {1,1} \rangle^+,&&\Ket | {1,2} \rangle^+,&&\Ket | {2,2} \rangle^+
\end{align}
zu finden. (6 Punkte)
\item[d)] Was ist die Summe der drei entsprechenden Wahrscheinlichkeiten? (2 Punkte)
\item[e)] Betrachten Sie nun einen $N$-Teilchenzustand $\Ket | {\underbrace{abb\dots b}_N} \rangle^+$. 
Was ist die Wahrscheinlichkeitsamplitude, bei einer Energiemessung den Zustand $\Ket | {\underbrace{1\dots 1}_N} \rangle^+$ zu erhalten? 
Diese Zustände seien ernaut auf 1 normiert. (6 Punkte)
\end{enumerate}
\newpage
\subsubsection*{Lösungsvorschlag}
\begin{enumerate}
\item[a)] Zunächst einmal gilt
\begin{align*}
\Ket | {a} \rangle&=c\Ket | {1} \rangle+s\Ket | {2} \rangle&&\tecxt\text{und}&\Ket | {b} \rangle&=-s\Ket | {1} \rangle+c\Ket | {2} \rangle.
\end{align*}
Damit lautet dann $\brak\langle {a}| {1}\rangle=c$, $\brak\langle {a}| {2}\rangle=s$, $\brak\langle {b}| {1}\rangle=-s$ und $\brak\langle {b}| {2}\rangle=c$.
Daraus folgen mit $\brak\langle {i,j}| {a,b}\rangle=\brak\langle {i}| {a}\rangle\brak\langle {j}| {b}\rangle$
\begin{align*}
\brak\langle {1,1}| {a,b}\rangle&=c(-s),&&&\brak\langle {1,2}| {a,b}\rangle&=cc,&&&\brak\langle {2,1}| {a,b}\rangle&=s(-s),&&&\brak\langle {2,2}| {a,b}\rangle&=sc.
\end{align*}
\item[b)] Die Summe der vier Wahrscheinlichkeiten ist
\begin{align*}
\sum_{ij}\vert\brak\langle {i,j}| {a,b}\rangle|^2=c^2s^2+c^2c^2+s^2s^2+c^2s^2=c^2(c^2+s^2)+s^2(c^2+s^2)=c^2+s^2=1.
\end{align*}
\item[c)] Für den symmetrisierten Zustand gilt $\Ket | {a,b} \rangle^+=\frac{1}{\sqrt{2}}\kla\left( {\Ket | {a,b} \rangle+\Ket | {b,a} \rangle} \right)$. 
Außerdem ist $\Ket | {1,2} \rangle^+=\frac{1}{\sqrt{2}}\kla\left( {\Ket | {1,2} \rangle+\Ket | {2,1} \rangle} \right)$.
Daraus folgen
\begin{align*}
{}^+\brak\langle {1,1}| {a,b}\rangle^+&=\frac{2}{\sqrt{2}}\brak\langle {1,1}| {a,b}\rangle=-\sqrt{2}cs,\\
{}^+\brak\langle {2,2}| {a,b}\rangle^+&=\frac{2}{\sqrt{2}}\brak\langle {2,2}| {a,b}\rangle=\sqrt{2}cs,\\
{}^+\brak\langle {1,2}| {a,b}\rangle^+&=\frac{2}{\sqrt{2}\sqrt{2}}\kla\left( {\brak\langle {1,2}| {a,b}\rangle+\brak\langle {2,1}| {a,b}\rangle} \right)=c^2-s^2.
\end{align*}
\item[d)] Die Summe der Wahrscheinlichkeiten ist
\begin{align*}
\sum_{ij}\left\vert {}^+\brak\langle {i,j}| {a,b}\rangle^+ \right\vert^2 =2c^2s^2+2c^2s^2+(c^2-s^2)^2=c^4+2c^2s^2+s^4=(c^2+s^2)^2=1.
\end{align*}
\item[e)] Für den symmetrisierten Zustand gilt explizit
\begin{align*}
\Ket | {\underbrace{abb\dots b}_N} \rangle&=\frac{1}{\sqrt{N}}\kla\left( {\Ket | {\underbrace{abb\dots b}_N} \rangle+\dots +\Ket | {\underbrace{bb\dots ba}_N} \rangle} \right).
\end{align*}
Daraus folgt für die Wahrscheinlichkeitsamplitude
\begin{align*}
{}^+\brak\langle {\underbrace{1\dots 1}_N}| {\underbrace{abb\dots b}_N}\rangle^+=\frac{N}{\sqrt{N}}\brak\langle {1}| {a}\rangle\brak\langle {1}| {b}\rangle^{N-1}=\sqrt{N}c(-s)^{N-1}.
\end{align*}
\end{enumerate}
\newpage
\section{Identische Teilchen (20 Punkte)}
Gegeben sei ein fermionisches Vielteilchensystem mit
fermionischen Erzeugungs- bzw. Vernichtungsoperatoren mit $\{c_i,c_j\}=0$ und $\{c_i,c_j^\dagger\}=\delta_{ij}$. 
Betrachten Sie nun die drei Operatoren
\begin{align}
A&=\frac{1}{2}\kla\left( {c_1^\dagger c_1-c_2^\dagger c_2} \right),&&&B&=\frac{1}{2}\kla\left( {c_1^\dagger c_2+c_2^\dagger c_1} \right),&&&C&=\frac{-\I}{2}\kla\left( {c_1^\dagger c_2-c_2^\dagger c_1} \right).
\end{align}
\begin{enumerate}
\item[a)] Berechnen Sie den Kommutator $[A, B]$.
Woran erinnert Sie die entstehende Vertauschungsrelation? (8 Punkte)
\item[b)] Der Vakuumzustand $|0\rangle $ ist durch die Eigenschaft $c_j |0\rangle = 0$ sowie $\langle 0|0\rangle = 1$ definiert.\\ 
Bestimmen Sie \hfill (3 Punkte)
\begin{align}
\braket\langle {0} | {A} | {0}\rangle,&&\braket\langle {0} | {B} | {0}\rangle,&&\braket\langle {0} | {C} | {0}\rangle.
\end{align}
\item[c)] Beweisen Sie die Relation $c_1c_1^\dagger\Ket | {0} \rangle=\Ket | {0} \rangle$. (2 Punkte)
\item[d)] In Besetzungszahldarstellung definieren wir $\Ket | {10} \rangle=c_1^\dagger\Ket | {0} \rangle$.\\
Bestimmen Sie \hfill (7 Punkte)
\begin{align}
\brak\langle {10}| {10}\rangle,&&\braket\langle {10} | {A} | {10}\rangle,&&\braket\langle {10} | {A^2+B^2+C^2} | {10}\rangle.
\end{align}
\end{enumerate}
\newpage
\subsubsection*{Lösungsvorschlag}
\begin{enumerate}
\item[a)] Definiere zunächst $n_j:=c_j^\dagger c_j$ und $a:=c_1^\dagger c_2$.
Dann gilt
\begin{align*}
A&=\frac{1}{2}\kla\left( {n_1-n_2} \right),&&&B&=\frac{1}{2}\kla\left( {a+a^\dagger} \right),&&&C&=\frac{-\I}{2}\kla\left( {a-a^\dagger} \right).
\end{align*}
Damit kann der Kommutator $[A,B]$ in einfachere Bestandteile zerlegt werden.
Für diese gilt
\begin{align*}
\klam\left[ {n_1,a} \right]&=c_1^\dagger c_1c_1^\dagger c_2-c_1^\dagger c_2c_1^\dagger c_1=c_1^\dagger \{c_1, c_1^\dagger\}c_2-c_1^{\dagger 2}c_1c_2+c_1^{\dagger 2}c_2c_1=c_1^\dagger c_2=a,\\
\klam\left[ {n_1,a^\dagger} \right]&=\klam\left[ {a,n_1^\dagger} \right]^\dagger=-a^\dagger,\\
\klam\left[ {n_2,a^\dagger} \right]&=\klam\left[ {n_1,a} \right]_{1\leftrightarrow 2}=(c_1^\dagger c_2)_{1\leftrightarrow 2}=c_2^\dagger c_1=a^\dagger,\\
\klam\left[ {n_2,a} \right]&=\klam\left[ {a^\dagger,n_2^\dagger} \right]^\dagger = -a,
\end{align*}
wobei $c_j^2=0=c_j^{\dagger 2}$ verwendet wurden. 
Daraus folgt dann 
\begin{align*}
\klam\left[ {A,B} \right]&=\frac{1}{4}\kla\left( {\klam\left[ {n_1,a} \right]+\klam\left[ {n_1,a^\dagger} \right]+\klam\left[ {-n_2,a} \right]+\klam\left[ {-n_2,a^\dagger} \right]} \right)=\\
&=\frac{1}{4}\kla\left( {a-a^\dagger +a-a^\dagger} \right)=\frac{1}{2}\kla\left( {a-a^\dagger} \right)=\I C.
\end{align*}
Dieser Zusammenhang erinnert an die Struktur einer Drehimpulsalgebra.
\item[b)] Aus $c_j |0\rangle = 0$ folgt auch $n_j\Ket | {0} \rangle=c_j^\dagger c_j\Ket | {0} \rangle=0$, $a\Ket | {0} \rangle=0$ und $a^\dagger\Ket | {0} \rangle=0$.
Somit gilt trivialerweise
\begin{align}
\braket\langle {0} | {A} | {0}\rangle&=0,&&&\braket\langle {0} | {B} | {0}\rangle&=0,&&&\braket\langle {0} | {C} | {0}\rangle&=0.
\end{align}
\item[c)] Es gilt $c_1c_1^\dagger = \{c_1,c_1^\dagger\}-c_1^\dagger c_1=1-n_1$ und somit $c_1c_1^\dagger\Ket | {0} \rangle=(1-n_1)\Ket | {0} \rangle=\Ket | {0} \rangle$.
\item[d)] Es gilt
\begin{align}
\brak\langle {10}| {10}\rangle&=\braket\langle {0} | {c_1c_1^\dagger} | {0}\rangle=\brak\langle {0}| {0}\rangle=1.
\end{align}
Für die anderen Ausdrücke ist es zunächst hilfreich drei Kommutatoren zu berechnen. 
Mit $\klam\left[ {\alpha\beta,\gamma} \right]=\alpha\{\beta,\gamma\}-\{\alpha,\gamma\}\beta$ lauten diese
\begin{align*}
\klam\left[ {A,c_1^\dagger} \right]&=\frac{1}{2}\kla\left( {\klam\left[ {c_1^\dagger c_1,c_1^\dagger} \right]-\klam\left[ {c_2^\dagger c_2,c_1^\dagger} \right]} \right)=\\
&=\frac{1}{2}\kla\left( {c_1^\dagger\{c_1,c_1^\dagger\}-\{c_1^\dagger,c_1^\dagger\}c_1-c_2^\dagger\{c_2,c_1^\dagger\}+\{c_2^\dagger,c_1^\dagger\}c_2} \right)=\frac{1}{2}c_1^\dagger,\\
\klam\left[ {B,c_1^\dagger} \right]&=\frac{1}{2}\kla\left( {\klam\left[ {c_1^\dagger c_2,c_1^\dagger} \right]+\klam\left[ {c_2^\dagger c_1,c_1^\dagger} \right]} \right)=\frac{1}{2}\kla\left( {0-0+c_2^\dagger\{c_1,c_1^\dagger\}-0} \right)=\frac{1}{2}c_2^\dagger,\\
\klam\left[ {C,c_1^\dagger} \right]&=\frac{-\I}{2}\kla\left( {\klam\left[ {c_1^\dagger c_2,c_1^\dagger} \right]-\klam\left[ {c_2^\dagger c_1,c_1^\dagger} \right]} \right)=\frac{-\I}{2}\kla\left( {0-0-c_2^\dagger +0} \right)=\frac{\I}{2}c_2^\dagger.
\end{align*}
Damit folgen dann
\begin{align*}
\braket\langle {10} | {A} | {10}\rangle&=\braket\langle {10} | {\klam\left[ {A,c_1^\dagger} \right]-c_1^\dagger A} | {0}\rangle=\frac{1}{2}\braket\langle {10} | {c_1^\dagger} | {0}\rangle=\brak\langle {10}| {10}\rangle=\frac{1}{2},\\
\braket\langle {10} | {A^2} | {10}\rangle&=\braket\langle {0} | {A\klam\left[ {A,c_1^\dagger} \right]-Ac_1^\dagger A} | {0}\rangle=\frac{1}{2}\braket\langle {10} | {A c_1^\dagger} | {0}\rangle=\frac{1}{2}\braket\langle {10} | {A} | {10}\rangle=\frac{1}{4},\\
\braket\langle {10} | {B^2} | {10}\rangle&=\braket\langle {10} | {B\klam\left[ {B,c_1^\dagger} \right]} | {0}\rangle=\frac{1}{2}\braket\langle {10} | {Bc_2^\dagger} | {0}\rangle=\\
&=\frac{1}{2}\braket\langle {10} | {\klam\left[ {B,c_1^\dagger} \right]_{1\leftrightarrow 2}} | {0}\rangle=\frac{1}{4}\braket\langle {10} | {(c_2^\dagger)_{1\leftrightarrow 2}} | {0}\rangle=\frac{1}{4},\\
\braket\langle {10} | {C^2} | {10}\rangle&=\braket\langle {10} | {C\klam\left[ {C,c_1^\dagger} \right]} | {0}\rangle=\frac{\I}{2}\braket\langle {10} | {Cc_2^\dagger} | {0}\rangle=\\
&=\frac{\I}{2}\braket\langle {10} | {\klam\left[ {-C,c_1^\dagger} \right]_{1\leftrightarrow 2}} | {0}\rangle=-\frac{\I^2}{4}\braket\langle {10} | {c_1^\dagger} | {0}\rangle=\frac{1}{4}.
\end{align*}
Aus den letzten drei Relationen folgt auch 
\begin{align*}
\braket\langle {10} | {A^2+B^2+C^2} | {10}\rangle=\frac{3}{4}\equiv S(S+1)
\end{align*}
mit $S=\frac{1}{2}$.
\end{enumerate}
\newpage
\section{Streuung und Born´sche Näherung (20 Punkte)}
Betrachten Sie die Streuung eines Teilchens mit einlaufender ebener Welle $\e^{\I kz}$ in $z$-Richtung an den Potentialen
\begin{align}
V_1(\vec{x})&=2V_0\delta^{(3)}(\vec{x})\\
V_2(\vec{x})&=V_0\kla\left( {\delta^{(3)}(\vec{x})+\delta^{(3)}(\vec{x}-\vec{a})} \right),
\end{align}
wobei $\vec{a}=(0,0,a)^\top$.
\begin{enumerate}
\item[a)] Geben Sie die Streuamplitude $f_1^{(1)}(\theta,\phi)$ für $V_1$ in erster Born'scher Näherung an. (4 Punkte)
\item[b)] Geben Sie die Streuamplitude $f_2^{(1)}(\theta,\phi)$ für $V_2$ in erster Born´scher Näherung an. 
Nutzen Sie im Ergebnis die Größe \hfill (5 Punkte)
\begin{align}
A&:=\vec{q}\cdot\vec{a}=ka(\cos\theta -1)=-2ka\sin^2\frac{\theta}{2}.
\end{align}
\item[c)] Skizzieren Sie das Schaubild von $\tecxt\text{d}\sigma / \tecxt\text{d}\Omega$ für beide Potentiale als Funktion von $\cos\theta$ für $ka=0.01$ und $ka=2\pi$. (6 Punkte)
\item[d)] Es sei $k$ fixiert und $|A|\ll 1$. 
Für welchen Winkel $\theta$ ist der Unterschied zwischen $(\tecxt\text{d}\sigma / \tecxt\text{d}\Omega)_{V_1}$ und $(\tecxt\text{d}\sigma / \tecxt\text{d}\Omega)_{V_2}$ maximal?
Wie groß muß $k$ mindestens sein, damit der relative Unterschied mehr als 1 Prozent beträgt (Näherungslösung reicht)? (5 Punkte)
\end{enumerate}
\newpage
\subsubsection*{Lösungsvorschlag}
\begin{enumerate}
\item[a)] Die Streuamplitude in erster Born'scher lautet
\begin{align*}
f_1^{(1)}(\theta,\phi)=-\frac{m}{2\pi}\intr\int_{\mathbb{R}^{3}}{V_1(\vec{x})\e^{-\I \vec{q}\cdot\vec{x}}}\,\mathrm{d}^{3}x=-\frac{V_0m}{\pi}.
\end{align*}
\item[b)] Die Streuamplitude in erster Born'scher lautet
\begin{align*}
f_2^{(1)}(\theta,\phi)=-\frac{m}{2\pi}\intr\int_{\mathbb{R}^{3}}{V_2(\vec{x})\e^{-\I \vec{q}\cdot\vec{x}}}\,\mathrm{d}^{3}x=-\frac{V_0m}{2\pi}-\frac{V_0m}{2\pi}\e^{-\I\vec{q}\cdot\vec{a}}=-\frac{\kappa}{2}\kla\left( {1+\e^{-\I A}} \right) .
\end{align*}
Dabei ist $\kappa:=\frac{V_0m}{\pi}$ und $A=ka(\cos\theta -1)$.
\item[c)] Der Plot in \ref{figure:DifferentialCrossSections} zeigt die differentiellen Wirkungsquerschnitte 
\begin{align*}
\kla\left( {\frac{\tecxt\text{d}\sigma}{\tecxt\text{d}\Omega}} \right)_{V_1}&=|-\kappa|^2=\kappa^2,\\
\kla\left( {\frac{\tecxt\text{d}\sigma}{\tecxt\text{d}\Omega}} \right)_{V_2}&=\frac{\kappa^2}{4}\kla\left( {1+\e^{-\I A}} \right)\kla\left( {1+\e^{\I A}} \right)=\frac{\kappa^2}{2}\kla\left( {1+\cos A} \right)
\end{align*}
für die geforderten Parameter.
Beachte wieder $A=ka(\cos\theta -1)$.
\begin{figure}[!h]
\centering
\includegraphics[width=\linewidth]{DifferentialCrossSections.png}
\caption{Differentielle Wirkungsquerschnitte für $ka\in\{0.1,2\pi\}$.
Der Wert $ka=0.1$ zeigt den Anstieg der Funktion, welcher für $ka=0.01$ kaum zu erkennen ist.}
\label{figure:DifferentialCrossSections}
\end{figure}
\item[d)] Betrachte die Differenz 
\begin{align*}
\Delta&:=\kla\left( {\frac{\tecxt\text{d}\sigma}{\tecxt\text{d}\Omega}} \right)_{V_1}-\kla\left( {\frac{\tecxt\text{d}\sigma}{\tecxt\text{d}\Omega}} \right)_{V_2}=\kappa^2-\frac{\kappa^2}{2}\kla\left( {1+\cos A} \right)\kom\overset{\substack{{A\,\ll\,1} \\ |}}{\simeq}\frac{\kappa^2}{4}A^2=\\
&=\kappa^2(ka)^2\sin^4\frac{\theta}{2}.
\end{align*}
Offensichtlich ist $\Delta$ maximal für $\theta=\pi$, also Rückstreuung. 
Hier ist dann $|A|=2ka$ und die relative Differenz beträgt
\begin{align*}
\frac{\Delta}{\kla\left( {\frac{\tecxt\text{d}\sigma}{\tecxt\text{d}\Omega}} \right)_{V_1}}&\simeq (ka)^2\sin^4\frac{\theta}{2}\Big\vert_{\theta=\pi}=(ka)^2.
\end{align*}
Damit diese Differenz größer als 1 Prozent ist, muss $ka>\frac{1}{10}$ gelten.
\end{enumerate}
\newpage
\section{Streuung und Partialwellen (20 Punkte)}
Es wird die Streuung eines Teilchens an einigen Potentialen mit der Partialwelle für $l = 0$ untersucht. 
Alle Potentiale haben die Reichweite $R_0$ und der betrachtete Impuls ist $k$, wobei $kR_0 = 14\pi$. 
Jemand hat bereits die radialen Schrödingergleichungen $(\partial_r^2 + k^2 - v(r))u(r) = 0$ für die Potentiale $v(r)$ exakt gelöst und die Ergebnisse für $u(r)$ als Funktion von $kr$ geplottet, siehe die beiden Plots.
\begin{enumerate}
\item[a)] Beschreiben Sie, welches Verhalten die radiale Wellenfunktion $u(r)$ im Plot aufweist und
warum. 
Beschreiben Sie, wie Sie aus dem Plot die Streuphase $\delta_0$ ablesen können. (5 Punkte)
\item[b)] Lesen Sie dann jeweils für beide Potentiale die Streuphase $\delta_0$ mit vernünftiger Genauigkeit ab und bestimmen Sie damit die Streuamplitude $f(\theta)$ in der Näherung,
dass nur $l = 0$-Streuung beiträgt. (10 Punkte) 
\item[c)] Welches der beiden Potentiale ist anziehend $(v(r) < 0)$, welches abstoßend $(v(r) > 0)$? (5 Punkte)
\end{enumerate}
\begin{figure}[h!]
\centering
\includegraphics[width=\linewidth]{u_for_v1.png}
\caption{$u(\rho=kr)$ für das Potential $v_1(r)$.}
\label{figure:uForV1}
\includegraphics[width=\linewidth]{u_for_v2.png}
\caption{$u(\rho=kr)$ für das Potential $v_2(r)$.}
\label{figure:uForV2}
\end{figure}
\newpage
\subsubsection*{Lösungsvorschlag}
\begin{enumerate}
\item[a)] Beschreiben Sie, welches Verhalten die radiale Wellenfunktion $u(r)$ im Plot aufweist und
warum. 
Beschreiben Sie, wie Sie aus dem Plot die Streuphase $\delta_0$ ablesen können. (5 Punkte)
\item[b)] Lesen Sie dann jeweils für beide Potentiale die Streuphase $\delta_0$ mit vernünftiger Genauigkeit ab und bestimmen Sie damit die Streuamplitude $f(\theta)$ in der Näherung,
dass nur $l = 0$-Streuung beiträgt. (10 Punkte) 
\item[c)] Welches der beiden Potentiale ist anziehend $(v(r) < 0)$, welches abstoßend $(v(r) > 0)$? (5 Punkte)
\end{enumerate}
%\begin{minted}[frame=lines, bgcolor=bg, fontsize=\footnotesize, linenos]{python}

%\begin{thebibliography}{9}
%\bibitem{example} description, \url{https://www.exmapleurl}, day.\,month~2021.
%\end{thebibliography}
\newpage
\appendix
\section{Nützliche Formeln und Ausdrücke}
Pauli-Matrizen
\begin{align}
\sigma^1&=\begin{pmatrix}
0&1\\ 1&0
\end{pmatrix},&&&\sigma^2&=\begin{pmatrix}
0&-\I\\ \I&0
\end{pmatrix},&&&\sigma^3&=\begin{pmatrix}
1&0\\ 0&-1
\end{pmatrix}.\label{eqn:PauliMatrizen}
\end{align}
$\gamma$-Matrizen in Diracdarstellung
\begin{align}
\gamma^0&=\begin{pmatrix}
I_2 & 0 \\ 0 & -I_2
\end{pmatrix},&&&\gamma^j&=\begin{pmatrix}
0&\sigma^j\\-\sigma^j &0
\end{pmatrix}. \label{eqn:GammaMatrizenDirac}
\end{align}
(Lorentz) Rotations-Matrizen für Diracspinoren
\begin{align}
S(R_x(\theta))&=\e^{-\I S_x\theta}=\begin{pmatrix}
\cos\frac{\theta}{2}&-\I\sin\frac{\theta}{2}&0&0\\
-\I\sin\frac{\theta}{2}&\cos\frac{\theta}{2}&0&0\\
0&0&\cos\frac{\theta}{2}&-\I\sin\frac{\theta}{2}\\
0&0&-\I\sin\frac{\theta}{2}&\cos\frac{\theta}{2}
\end{pmatrix}, \label{eqn:RotationsMatrizenSX}\\
S(R_y(\theta))&=\e^{-\I S_y\theta}=\begin{pmatrix}
\cos\frac{\theta}{2}&-\sin\frac{\theta}{2}&0&0\\
\sin\frac{\theta}{2}&\cos\frac{\theta}{2}&0&0\\
0&0&\cos\frac{\theta}{2}&-\sin\frac{\theta}{2}\\
0&0&\sin\frac{\theta}{2}&\cos\frac{\theta}{2}
\end{pmatrix}, \label{eqn:RotationsMatrizenSY}\\
S(R_z(\theta))&=\e^{-\I S_z\theta}=\begin{pmatrix}
\e^{-\I\frac{\theta}{2}}&0&0&0\\
0&\e^{\I\frac{\theta}{2}}&0&0\\
0&0&\e^{-\I\frac{\theta}{2}}&0\\
0&0&0&\e^{\I\frac{\theta}{2}}
\end{pmatrix}. \label{eqn:RotationsMatrizenSZ}
\end{align}

\end{document}